\documentclass[11pt]{article}
\usepackage{microtype,amsmath,amssymb,amsthm,mathtools,booktabs,array}
\usepackage[margin=1.05in]{geometry}
\usepackage[colorlinks=true,linkcolor=blue!45!black,citecolor=blue!45!black,urlcolor=blue!45!black]{hyperref}
\usepackage{xcolor}
\hypersetup{pdftitle={The dimension-degree boundary for Gaussian moment vanishing},pdfsubject={Working mathematical manuscript}}
\newtheorem{theorem}{Theorem}[section]
\newtheorem{proposition}[theorem]{Proposition}
\newtheorem{lemma}[theorem]{Lemma}
\newtheorem{corollary}[theorem]{Corollary}
\theoremstyle{definition}\newtheorem{definition}[theorem]{Definition}
\theoremstyle{remark}\newtheorem{remark}[theorem]{Remark}
\newcommand{\E}{\mathbb E}
\newcommand{\C}{\mathbb C}
\newcommand{\tr}{\operatorname{tr}}
\newcommand{\wt}{\operatorname{wt}}
\newcommand{\G}{\mathsf G}
\DeclareMathOperator{\Span}{span}
\title{Uniform Gaussian cut decay and the ordinary-node center axes}
\author{}\date{17 November 2026: internal written proof}
\begin{document}\maketitle
This note proves a sufficient criterion in every dimension and polynomial
degree,then completes the three-variable ordinary-node center-axis family.
The analytic proof retains complete actual cut homology,uniform tails,
and mass and boundary-mass control before taking a current limit. Its
classical ingredients are explicitly credited. The criterion has a
uniform bound for the FULL phase gradient;bounded critical points alone
do not suffice. No novelty,independent replay,formal verification or
external review is claimed.

For independent standard Gaussian variables,after any invertible affine
coordinate change retain the full nondegenerate affine quadratic $q$ and
its normalized formal expectation $\mathbb E$. Set
\[
 \Phi_A(t)=\sum_{m\ge0}\mathbb E[AP^m]\,t^m/m!,\qquad
 M=\mathbb C(t)[x]/\sum_i(\partial_i+tP_i-q_i)\mathbb C(t)[x].
\]
The connection is $\nabla_t=\partial_t+P$. Gaussian integration by
parts defines a horizontal functional $M\to\mathbb C((t))$.
Purity means $\Phi_1=1$;safety means every $\Phi_A$ is polynomial.
The quotient is an ambient module;no image-kernel computation is needed.
\begin{lemma}[The seed functional]\label{lem:seedfunctional}
Let $M$ be a finite-dimensional differential module over $K=\C(t)$,
let $s\in M$, and suppose a $K$-linear horizontal map
$J:M\longrightarrow\C((t))$ satisfies $J(s)=1$. Then the differential
submodule $U=K\langle\partial_t\rangle s$ has a horizontal quotient
$L:U\longrightarrow K$ with $L(s)=1$.

At an ordinary base point, $L$ extends locally to an analytic horizontal
functional on $M$. For the period conclusion, assume explicitly that
$M$ is the generic exponential de Rham system of regular functions
$f,g$ on a smooth affine complex variety, where $f$ is the base
projection and $f,g$ are algebraically independent. Remove the finite
exceptional base set in the Hien--Roucairol comparison theorem.
Then the extension is represented by a flat rapid-decay cycle, and its
period evaluated on $s$ is identically one.
\end{lemma}
\begin{proof}
On the generators $s,\nabla s,\nabla^2s,\ldots$, the values of $J$
are $1,0,0,\ldots$. Therefore its restriction takes values in $K$.
Every rational relation between these generators is respected because
$J$ was defined on $M$, so this restriction is well defined and
horizontal. Its value at $s$ makes it surjective.

Choose a point where $M$, $U$ and $L$ are ordinary, extend the fiber
functional from $U$ to $M$, and solve the dual analytic ODE.
Uniqueness preserves its restriction to $U$. The perfect period
pairing and solution representation in
Hien--Roucairol~\cite[Theorems~2.4 and~3.7]{HienRoucairol} give the
asserted cycle. A scalar rational normalization is an invertible gauge
away from its zeros and poles and must be retained in the period.
\end{proof}
\begin{lemma}[A positive original seed]\label{lem:positiveformal}
Suppose a finite differential module $M$ has,after a finite cover
$r=t^{-1/N}$,a full frame of rank-one formal exponential factors.
If every zero-exponential horizontal functional takes a seed $a$ into
$r\mathbb C[[r]]$,the formal Gaussian functional cannot take $a$ to one.
\end{lemma}
\begin{proof}
If its value on $a$ is one,its restriction to the rational differential
submodule generated by $a$ takes values in $\mathbb C(t)$:the generator
values are $1,0,0,\ldots$. Thus it defines a horizontal quotient to the
trivial module. Formal extension is exact. After ramification the full
module is a direct sum of rank-one factors,so this quotient extends from
the seed submodule to the full module. A nonzero exponential factor
has no horizontal map to the trivial module over $\mathbb C((r))$:
the logarithmic derivative of a Laurent series cannot have its higher
polar derivative. On the zero part all horizontal maps are constant
linear combinations of the frame. Their seed values have positive order
and cannot sum to one.
\end{proof}
\subsection{A uniform cut criterion in any degree}
\label{sec:uniformcutdecay}

Intersecting critical curves need not admit a simple formal normal
form. The following criterion avoids choosing one. Its hypothesis
concerns the full phase $tP-q$, uniformly in large complex $t$;
bounded critical points alone would not suffice.

\begin{theorem}\label{thm:uniformcuts}
Let $P,q\in\C[x_1,\ldots,x_n]$, with $P$ nonconstant. Suppose
there are fixed $c,R_0,t_0>0$ such that
\begin{equation}\label{eq:uniformphasegradient}
 \|\nabla_x(tP-q)(x)\|\ge c\|x\|
       \quad(|t|\ge t_0,\ \|x\|\ge R_0).
\end{equation}
Along every fixed ray $t=Re^{i\theta}$, the full rapid-decay
homology has a flat cut basis $\Gamma_i(t)$, attached to the
critical values $c_i(t)$ of $tP-q$, such that for every fixed
polynomial $a$,
\begin{equation}\label{eq:normalizedcutdecay}
 e^{-c_i(t)}\int_{\Gamma_i(t)}a(x)e^{tP-q}\,
                 dx_1\cdots dx_n\longrightarrow0.
\end{equation}
Each $c_i(t)=\alpha_i t+O(1)$, where $\alpha_i$ is a critical
value of $P$. If all critical values of $P$ lie in a closed
half-plane through zero, a supporting $t$ ray makes every flat
period of every such fixed form tend to zero.
\end{theorem}

We give the completeness, tail and limiting-form arguments separately.
None follows just from counting the phase-critical points.

\begin{proof}
For every sufficiently large $t$, \eqref{eq:uniformphasegradient}
makes the phase-critical set finite and contained in one fixed
ball. It also gives complex tameness at infinity. Pham's tame
Milnor decomposition~\cite[Part~I and Appendix]{Pham} identifies
the complete middle rapid-decay homology with the local vanishing
homologies carried along cuts from its critical values. The rank
is the sum of their Milnor numbers, including degenerate points.
If this rank is zero the conclusion is immediate. Otherwise
retain all these local generators; identical critical values
are grouped, not discarded.

On the chosen ray, the critical points and values are algebraic
functions of $R$, after a finite cover. Enlarge $R_1$ so branches
and multiplicities stabilize. Choose $|\eta|=1$ with
$\Re\eta\ge1/2$ so the cuts
\[
 S_t(x)=c_i(t)-\eta s,\qquad s\ge0,\quad S_t=tP-q,
\]
are disjoint for $R>R_1$, except at identical grouped endpoints.
There are only finitely many forbidden choices: every nonzero
difference of two algebraic critical values has one limiting
direction. The cut angle is independent of $\theta$. In
particular a supporting $t$ ray need not be a generic ray.

Along a cut, normalized complex gradient transport is
\[
 x_s=-\eta\,\overline{\nabla S_t}/\|\nabla S_t\|^2,
               \qquad S_t(x(s))_s=-\eta.
\]
Outside the fixed ball,
$|d\|x\|^2/ds|\le2/c$. All local vanishing generators start
in that ball. With fixed sufficiently large $A$ and $B>2/c$,
they therefore have representatives inside the carrier
\begin{equation}\label{eq:uniformcarrier}
 \|x\|^2\le A+Bs.
\end{equation}
At its outer boundary the transport points strictly inward
relative to this inequality. No trajectory escapes at finite
cut length, and the cuts avoid all other critical values.

We must control a single family of representatives. The gradient
flow itself need not be semialgebraic. Instead compactify $x$
by $x/\sqrt{1+\|x\|^2}$ and $s$ by $s/(1+s)$. Take the closure
of each constrained cut carrier and mark its rapid boundary
$G$, where $s=\infty$. Inequality~\eqref{eq:uniformcarrier}
forbids spatial infinity at finite $s$. For fixed $R$, the
relative chains of this compact pair are precisely locally
finite cut chains with their ends allowed only at rapid
infinity. This follows also by a compatible triangulation of
the compact pair: its relative simplices restrict to the usual
locally finite chain model in the open carrier. The exponential
there is exactly $e^{c_i-\eta s}$.

For the full unconstrained cut carrier, mark and remove any
spatial boundary at finite $s$; it is not a rapid boundary.
The constrained carrier, the full carrier, their rapid boundaries
and this excluded boundary form a finite family of semialgebraic
subsets in the same compactification. Apply the compatible
finite-family form of Hardt triviality to projection onto $R$.
It gives one trivialization on an unbounded interval, preserving
these subsets and their inclusions. The precise form is stated
in Drusvyatskiy--Lewis~\cite[Theorem~2.25]{DrusvyatskiyLewis};
it is Hardt's theorem. Semialgebraic triangulation and the
chain-current boundary rule are as in
Hardt--Lambrechts--Turchin--Voli\'c~\cite[Sections~2.1 and~3,
especially Theorem~3.5]{HLTV}.

Gradient transport already showed that the constrained inclusion
contains every local vanishing generator. At one parameter choose
finitely many integral relative chains whose images are a basis
over $\C$ of the full rapid-decay space. A nonprimitive integral
image causes no problem over $\C$. Transport these chains by
the compatible trivialization. They form a jointly semialgebraic
family with fixed integer multiplicities and supports satisfying
\eqref{eq:uniformcarrier}. Their images are a flat basis of the
actual homology local system: on each compact parameter interval
the transport preserves the relative pairs and is proper on
every finite phase band. Thus this is a completeness statement
for actual cycles, not an identification of presumed formal blocks.

Here is the needed uniform mass bound. In a fixed semialgebraic
family of $d$-chains with bounded multiplicities, the mass inside
a ball of radius $L\ge1$ is at most $C L^d$, uniformly in the
parameters. Partition tangent planes by their largest coordinate
$d$-minor. The Cauchy--Binet identity gives a lower bound
$\binom Nd^{-1/2}$ for some such minor in ambient dimension $N$.
The area formula then bounds that piece by its coordinate
projection volume times the number of preimages. Semialgebraic
uniform finiteness bounds those finite fibers independently of
the parameters; see, for example,
\cite[Theorem~2.32]{DrusvyatskiyLewis}. Sum the finitely many
pieces and multiplicities. Apply the same argument to boundary
chains. This proves uniformity, rather than merely pointwise
finite volume.

Lift the cut chains to $(x,s)$. Their bands $j\le s\le j+1$
lie in balls of radius $O(1+j)$ by
\eqref{eq:uniformcarrier}. Their fixed semialgebraic descriptions
and the mass estimate give, with constants independent of $R$,
\begin{equation}\label{eq:uniformcuttail}
 \left|e^{-c_i(t)}\int_{s\ge L}a e^{S_t}\,dx_1\cdots dx_n\right|
 \le C\sum_{j\ge\lfloor L\rfloor}(1+j)^M e^{-j/2}.
\end{equation}
The polynomial amplitude changes only $M$. This tail tends to
zero uniformly in $R$.

It remains to bound a fixed phase band. A bounded critical branch
$x(r)$, $r=1/t$, satisfies
\[
 q(x(r))=-c_i(r)-r c_i'(r).
\]
At a stationary point the chain rule gives $c_i'(r)=-P(x(r))/r^2$;
eliminating $P$ gives this identity. Since $q(x(r))$ is
bounded, its Puiseux expansion forces the only possible polar
term of $c_i$ to be $\alpha_i/r$. Its coefficient is
$\alpha_i=P(x(0))$, and $\nabla P(x(0))=0$. Hence
$c_i(t)=\alpha_i t+O(1)$ even at critical collisions.

Choose a generic fixed truncation level $L$. The truncated lifted
chains have compact supports and uniformly bounded mass and
boundary mass. The latter follows from the same semialgebraic
estimate applied to their slices; after a finite stratification
there is an arbitrarily large sequence of such levels. By the
Federer--Fleming compactness theorem, every sequence $R\to\infty$
has a subsequence converging to an integral $n$-current. We use
the version with both mass and boundary-mass bounds,
Simon~\cite[Chapter~6, Section~3, Theorem~3.15]{SimonGMT}.
Its rectifiability conclusion is essential.

On the lifted chain,
\[
 P(x)-\alpha_i=
       \frac{c_i(t)-\alpha_i t+q(x)-\eta s}{t}.
\]
Every limiting current is therefore supported in
$\{P=\alpha_i\}\times[0,L]$. Stratify this complex algebraic
level set into smooth complex strata. Almost everywhere for the
rectifiable limiting current, its tangent plane projects into
the tangent of one of these strata. Each has complex dimension
less than $n$. The holomorphic $n$-form $dx_1\cdots dx_n$ thus
vanishes on that tangent plane. Lower-dimensional strata carry
no $n$-mass; the same argument covers singular strata of larger
real dimension. Testing with the fixed smooth form
$e^{-\eta s}a(x)dx_1\cdots dx_n$, with a compact cutoff, gives
zero on every subsequential limit. Consequently the normalized
integral on $0\le s\le L$ tends to zero. This does not assert
that the real volume of the chains tends to zero.

First let $R\to\infty$, then let $L\to\infty$ through the
allowed levels in \eqref{eq:uniformcuttail}. This proves
\eqref{eq:normalizedcutdecay}. Finally choose $\theta$ with
$\Re(e^{i\theta}\alpha)\le0$ for all critical values $\alpha$
of $P$. The factors $e^{c_i(t)}$ are then bounded, so every basis
period tends to zero. A flat cycle is a constant linear
combination of this full basis, proving the last assertion.
\end{proof}

\begin{corollary}\label{cor:twocriticalvalues}
Let $q$ be a full affine Gaussian quadratic and assume
\eqref{eq:uniformphasegradient}. If $P$ has at most two critical
values, every original marked Gaussian period is either zero or
is not a finite exponential-polynomial. In particular $P$ is
neither pure nor safe.
\end{corollary}
\begin{proof}
Any two complex numbers lie in a closed half-plane through zero.
For two nonzero values normalize them to $z_1,z_2$ on the unit
circle. If $z_1+z_2\ne0$, use the direction
$-\overline{z_1+z_2}/|z_1+z_2|$; if their sum is zero, use
$i\overline{z_1}$. Zero or missing values add no condition.

A nonzero constant original marked period would, by the rational
seed restriction and actual-period extension in
Lemma~\ref{lem:seedfunctional}, be represented by a flat period
of that fixed polynomial form. Its generic exponential system
has $f=t$, $g=tP-q$ on smooth affine space and regular input
$\mathcal O$; $f,g$ are algebraically independent because $P$
is nonconstant. Remove the finite exceptional parameter set of
the Hien--Roucairol comparison. Theorem~\ref{thm:uniformcuts}
makes the resulting constant period tend to zero, a contradiction.
No convergence of the formal moment series was assumed.

A constant-coefficient operator reduces a nonzero finite
exponential-polynomial to $c e^{\beta t}$ while replacing its
original marker by another fixed polynomial in $P$ times that
marker. Shift $P$ by $-\beta$ and normalize. Its gradient is
unchanged and its critical values are still at most two. The
same contradiction applies. Since $\Phi_1(0)=1$, it cannot be
a polynomial, proving also the safety exclusion.
\end{proof}

\subsection{Completing the ordinary-node center axes}
\label{sec:ordinaryaxes}

The only remaining ordinary-node row is
\begin{equation}\label{eq:ordinaryaxes}
 P=Tuv+a u^3+d v^3+b u^2+c v^2+k.
\end{equation}
Retain the full Gaussian
\[ q=\tfrac12(Au^2+Bv^2+mT^2+2Cuv+2DuT+2EvT)+\ell_u u+\ell_v v+\ell_TT+q_0,\]
with nondegenerate Hessian and $m\ne0$. Assume also the following
conditions at the other leading nodes:
\[
 a=b=0\ \Longrightarrow\ A\ne0,\qquad
 d=c=0\ \Longrightarrow\ B\ne0.
\]
We use a full formal frame when $bc\ne0$ and the preceding
uniform cut criterion when $bc=0$. These branches are disjoint
and include every zero face of the four coefficients.

\begin{lemma}\label{lem:axisuniformgradient}
For \eqref{eq:ordinaryaxes} with these Gaussian conditions,
\eqref{eq:uniformphasegradient} holds.
\end{lemma}
\begin{proof}
Otherwise take $\|x\|\to\infty$, $|t|\ge t_0$ and
$\nabla(tP-q)=o(\|x\|)$. Division by $t\|x\|^2$ shows that
a limiting projective direction is a leading node: the $T$ axis,
also the $u$ axis if $a=0$, and the $v$ axis if $d=0$.

Near any such node use height $z$ and transverse coordinates
$x,y=o(z)$. The phase polynomial has the form
\[
 P=zxy+\beta z^2+H(x,y)+k,
\]
where $H$ has only terms of degrees two and three. There are
no transverse linear terms. The transverse equations, divided
by $tz$, give
\[
 |x|+|y|\le C/|t|+o(|x|+|y|)+o(1/|t|).
\]
The terms of $\nabla H/z$ are absorbed because $(x,y)=o(z)$;
the Gaussian cross terms divided by $tz$ are bounded by the
same quantities. Thus $x,y=O(1/|t|)$. The height equation is
\[
 \partial_z(tP-q)=(2\beta t-M)z+O(1),
\]
where $M$ is the Gaussian square coefficient along this node.
If $\beta\ne0$, $|2\beta t-M|$ is uniformly bounded below
for all sufficiently large $|t|$. If $\beta=0$, the no-frame
condition gives $M\ne0$. This contradicts the assumed small
gradient ratio. The proof also covers bounded-$t$ subsequences
within the chosen large-$|t|$ domain and therefore proves the
required uniform estimate.
\end{proof}

\begin{lemma}\label{lem:axisrank}
The full generic twisted rank of \eqref{eq:ordinaryaxes} is
\[
 N=5+\mathbf1_{a\ne0}+\mathbf1_{d\ne0}.
\]
\end{lemma}
\begin{proof}
The three homogenized quadratic gradient equations have total
intersection length eight in $\mathbb P^3$. Their points at
infinity are exactly the leading ordinary nodes, in number
$3-\mathbf1_{a\ne0}-\mathbf1_{d\ne0}$. Each has length one.
Indeed its local Jacobian has the two transverse columns of
$\operatorname{Hess}F_3(w)$ and the column $\nabla(tF_2-q_2)(w)$.
The first two have rank two and are annihilated by pairing with
$w$; the last has nonzero pairing
$2(tF_2-q_2)(w)$. This value is nonzero for large generic $t$
by the stated conditions. Thus these infinity intersections are
simple. The uniform gradient estimate gives a finite affine
critical set; there are no other components. B\'ezout and
Pham's tame comparison give the asserted full rank.
\end{proof}

\begin{proposition}\label{prop:axisexclusion}
Every phase \eqref{eq:ordinaryaxes} under the stated conditions
is not pure. Every original marked Gaussian period is zero or
is not a finite exponential-polynomial.
\end{proposition}
\begin{proof}
First let $bc\ne0$. Set $r=t^{-1/2}$, $u=rU$, $v=rV$, leaving
$T$ unchanged. Exactly,
\begin{align*}
 tP-q={}&k/r^2+F_0(U,V,T)+E_r(U,V,T),\\
 F_0={}&TUV+bU^2+cV^2-mT^2/2-\ell_TT-q_0,\\
 E_r={}&r(aU^3+dV^3-DUT-EVT-\ell_uU-\ell_vV)\\
       &-r^2(AU^2+BV^2+2CUV)/2,\\
 du\,dv\,dT={}&r^2dU\,dV\,dT.
\end{align*}
The full polynomial twisted quotient for $F_0$ has rank five,
including its degenerate parameter values. Here is a separate rank
proof. Suppose $\nabla F_0=o(\|(U,V,T)\|)$ on an unbounded
sequence. If $T$ stays bounded, the first two equations force
both normalized $U,V$ to be nonzero; the last, divided by the
squared norm, gives their product zero, a contradiction. If
$T$ is unbounded, the first two equations have inverse
$O(1/|T|)$, giving $U,V=o(1)$ and $\|(U,V,T)\|\asymp|T|$.
The last equation then has leading term $-mT$, again a
contradiction. This coercivity is uniform on compact sets with
$bcm\ne0$. The coefficient domain is connected, so conservation
of critical intersection multiplicity and Pham's tame comparison
make its rank constant. At generic $\ell_T$, its critical
equations give one point with $U=V=0$ and four with
$T^2=4bc$, all simple. Hence the rank is five, also when points
merge at the exceptional values of $\ell_T$.

Choose any polynomial basis $B_1,\ldots,B_5$ of this fixed
quotient and its dual $\lambda_1,\ldots,\lambda_5$. Then
\[
 L_j(a_0)=e^{kt}r^2\lambda_j\bigl(e^{E_r}a_0(rU,rV,T)\bigr)
\]
defines five formal horizontal functionals on the original full
module. Coefficients are finite polynomial evaluations at every
order. Spatial exactness is the fixed-phase conjugacy; time
horizontality retains the displayed $r^2$ volume and the
coordinate velocity $(u/r,v/r,0)$, whose divergence is $2/r$.
Cartan's formula makes the remaining velocity term spatially exact.
For every original fixed polynomial marker their values lie in
$e^{kt}r^2\C[[r]]$.

The other $P$-critical points are the off-axis points
$p_u=(-2b/(3a),0,0)$ if $a\ne0$ and
$p_v=(0,-2c/(3d),0)$ if $d\ne0$. They are Morse, with Hessian
determinants $8b^3/(9a^2)$ and $8c^3/(9d^2)$. Holomorphic Morse coordinates reduce their local phases to a
nondegenerate quadratic. Expanding the analytic Gaussian amplitude
and volume unit and applying the quadratic Wick recurrence gives
all-order horizontal characters. These have prefactors $e^{P(p)t}r^3$, multiplied by power series in
$r^2$ on fixed markers. Their values may coincide with each
other or with $k$; no separation hypothesis is used.

Here is an explicit completeness test. Let $M$ bound the degree
of the chosen $B_j$ in $U,V$. Substitute $U=u/r,V=v/r$ and
multiply each $B_j$ by
\[
 C(u,v)=\prod_{p_u}(1-u/u_p)^K
                      \prod_{p_v}(1-v/v_p)^K,\qquad K>M.
\]
An empty product is one. At the central fixed model this is
$1+O(r)$, so the normalized five-by-five matrix is the identity
at leading order. At each off-axis point its zero of order $K$
overcomes the worst $r^{-M}$ growth of the substituted basis;
the normalized cross values tend to zero. For the off-axis rows
use $(u/u_p)^K$ or $(v/v_p)^K$. They have zero central leading
values, unit value at their own Morse point and zero leading
value at the other. The resulting full matrix has invertible
block diagonal leading term. Its size is exactly the rank of
Lemma~\ref{lem:axisrank}. These are therefore a full formal
frame, including all fixed-model degeneracies and equal actions.

Every zero-exponential original-marker seed has positive order,
either $r^2$ or $r^3$. Lemma~\ref{lem:positiveformal} excludes
a nonzero constant original marked period. Differential
isolation and shifting $k$ then exclude every nonzero finite
exponential-polynomial as before.

It remains to treat $bc=0$. The critical equations of $P$ show
that its values are among
\[
 k,\qquad k+4b^3/(27a^2)\quad(ab\ne0),\qquad
 k+4c^3/(27d^2)\quad(dc\ne0).
\]
At a missing isolated point omit its displayed value. Any extra
critical axes have value $k$. Since $bc=0$, there are at most
two values, including every intersecting-axis or degenerate
case. Lemma~\ref{lem:axisuniformgradient} and
Corollary~\ref{cor:twocriticalvalues} apply. This completes both
disjoint branches and proves the proposition.
\end{proof}

\begin{thebibliography}{99}
\bibitem{DrusvyatskiyLewis}
D.~Drusvyatskiy and A.~S.~Lewis, \emph{Semi-algebraic functions have
small subdifferentials}, \emph{Mathematical Programming, Series~B}
\textbf{140} (2013), 5--29, Theorems~2.25 and~2.32;
\href{https://doi.org/10.1007/s10107-012-0624-x}{doi:10.1007/s10107-012-0624-x}.

\bibitem{HLTV}
R.~M.~Hardt, P.~Lambrechts, V.~Turchin and I.~Voli\'c,
\emph{Real homotopy theory of semi-algebraic sets},
\href{https://arxiv.org/abs/0806.0476v3}{arXiv:0806.0476v3},
Sections~2.1 and~3, especially Theorem~3.5.

\bibitem{SimonGMT}
L.~Simon, \emph{Introduction to Geometric Measure Theory},
author's 2017/2018 notes, Chapter~6, Section~3, Theorem~3.15,
printed page~175;
\href{https://math.stanford.edu/~lms/ntu-gmt-text.pdf}{author's primary text}.

\bibitem{HienRoucairol}
M.~Hien and C.~Roucairol, \emph{Integral representations for solutions
of exponential Gau\ss--Manin systems},
\emph{Bulletin de la Soci\'et\'e Math\'ematique de France}
\textbf{136} (2008), no.~4, 505--532,
\href{https://doi.org/10.24033/bsmf.2564}{doi:10.24033/bsmf.2564}.

\bibitem{Pham}
F.~Pham, \emph{La descente des cols par les onglets de Lefschetz,
avec vues sur Gauss--Manin}, \emph{Ast\'erisque} \textbf{130}
(1985), 11--47,
\href{https://numdam.org/item/AST_1985__130__11_0/}{Numdam primary text}.

\end{thebibliography}
\end{document}
